\documentclass[11pt,aps,preprint,nofootinbib,superscriptaddress,tightenlines]{revtex4-2}
\usepackage{amsmath,amssymb,bm,mathtools,graphicx}
\usepackage{xcolor}
\definecolor{darkblue}{RGB}{0,0,160}
\usepackage[colorlinks=true,linkcolor=darkblue,citecolor=darkblue,urlcolor=darkblue]{hyperref}

\def\st{\begin{equation}}
\def\stp{\end{equation}}
\def\Eq#1{eq.~(\ref{#1})}
\def\Sect#1{Sect.~\ref{#1}}
\def\Fig#1{Fig.~\ref{#1}}
\def\one{\mathbb 1}
\def\tr{\operatorname{tr}}
\let\oldparagraph\paragraph
\renewcommand{\paragraph}{\par\bigskip\oldparagraph}

\begin{document}

\title{One kernel for the ideal and the viscous step}
\author{D.~Teaney}
\affiliation{Department of Physics and Astronomy, Stony Brook University,
Stony Brook, New York 11794, USA}
\date{\today}

\begin{abstract}
The ideal step smooths with the SPH kernel $W(|\bm r|)$, round in the lab.  The viscous step weights its
bonds with $\psi(\rho)$ at the comoving distance $\rho^2=\bm r\cdot\ell^{-1}\bm r$, round in the rest frame of the
particle, so that the bond moment is $n_a\ell_a$.  We replace $\psi$ by the bond weight of the SPH kernel
itself, $\psi_{ab}=-W'(r_{ab})/r_{ab}$ at the lab distance, so that the bond moment is $n_a\one$, and put the
comoving metric $\ell_a$ in by hand.  For the stochastic step this changes only the weights and the
faces: the metric already sits in the Onsager matrix of the stress move.  For the deterministic step the
bond form is written in its exact resolved-plus-residual form.  This note works out the formulas,
describes the implementation (\texttt{SPH\_DISSIPATION=common}) and its tests.
\end{abstract}

\maketitle

%=============================================================================
\section{The present scheme}
\label{sec:present}

Notation as in \texttt{particle\_metrics.tex} and \texttt{stochastic.tex}: $\bm r_{ab}=\bm x_b-\bm x_a$, $r_{ab}=|\bm r_{ab}|$,
$V_a$ the particle volume, $\beta_a=\bm v_a/T_a$, $\Delta\beta_{ab}=\beta_b-\beta_a$, and the comoving metric
\st
\ell_a=\one-\bm v_a\bm v_a^T,\qquad \ell_a^{-1}=\one+\gamma_a^2\bm v_a\bm v_a^T .
\stp

\paragraph{The ideal step} uses the SPH kernel $W(r;h)$ in the lab, through $\bm D_{ab}=\partial W(\bm x_a-\bm x_b)/\partial\bm x_a$
and the number density $n_a=\sum_bW_{ab}$.

\paragraph{The viscous step} uses a bond weight $\psi_{ab}=\psi(\rho_{ab})$, $\psi(\rho)\propto(1-\rho/2h_\psi)^3$, evaluated at the
comoving distance
\st
\rho_{ab}^2=\bm r_{ab}\cdot\ell_a^{-1}\bm r_{ab} ,
\stp
an ellipse contracted by $1/\gamma$ along $\bm v_a$, corrected per particle by $\Lambda_a$ so that the bond moment is
exactly
\st
M_a\equiv\sum_{b\neq a}V_b\,\psi_{ab}\,\bm r_{ab}\bm r_{ab}^T=n_a\,\ell_a .
\label{Mold}
\stp
With the normalized weights $w_{ab}=V_b\psi_{ab}/n_a$, $\sum_bw_{ab}\bm r_{ab}\bm r_{ab}^T=\ell_a$, and the faces
\st
\bm A_{ab}=V_a\,w_{ab}\,\ell_a^{-1}\bm r_{ab},\qquad \sum_{b\neq a}\bm A_{ab}\,\bm r_{ab}^T=V_a\one .
\label{Aold}
\stp
Two kernels: $W$ round in the lab for the ideal step, $\psi$ round in the rest frame for the viscous step.

%=============================================================================
\section{One kernel}
\label{sec:common}

\paragraph{The weights.}  Take the bond weight from the SPH kernel itself, at the lab distance:
\st
\psi_{ab}=-\frac{W'(r_{ab};h)}{r_{ab}},\qquad\text{i.e.}\qquad \psi_{ab}\,\bm r_{ab}=\bm D_{ab}.
\label{psicommon}
\stp
The bond vector weighted by $\psi$ is the kernel gradient of the ideal step.  For the Wendland C2 kernel
$\psi\propto(1-r/2h)^3$, the present shape with $h_\psi=h$; for C4 and C6 it is the gradient of those kernels.
The weights are round in the lab, so the bond moment is
\st
M_a=\sum_{b\neq a}V_b\,\psi_{ab}\,\bm r_{ab}\bm r_{ab}^T=n_a\one
\label{Mnew}
\stp
up to sampling corrections, which the same $\Lambda_a$ removes (with the flat metric in place of $\ell_a$).
Then
\st
w_{ab}=\frac{V_b\psi_{ab}}{n_a},\qquad \sum_{b\neq a}w_{ab}\,\bm r_{ab}\bm r_{ab}^T=\one,\qquad
\bm A_{ab}=V_a\,w_{ab}\,\bm r_{ab},\qquad \sum_{b\neq a}\bm A_{ab}\,\bm r_{ab}^T=V_a\one .
\label{wnew}
\stp
At rest $\ell_a=\one$ and nothing changes: for C2 the scheme is the present one.

\paragraph{The fit} needs no metric (\texttt{particle\_metrics.tex}, Operation 1):
\st
G_a=\sum_{b\neq a}w_{ab}\,\Delta\beta_{ab}\,\bm r_{ab}^T\,(M_a/n_a)^{-1}=\frac1{V_a}\sum_{b\neq a}\Delta\beta_{ab}\,\bm A_{ab}^T ,
\qquad (G_a)_{ij}\simeq\partial_j\beta_i ,
\label{fit}
\stp
exact for linear fields.  With \Eq{psicommon} it is the moment-corrected SPH gradient.  The residuals
$\delta_{ab}=\Delta\beta_{ab}-G_a\bm r_{ab}$ are orthogonal to the bonds,
\st
\sum_{b\neq a}w_{ab}\,\delta_{ab}\,\bm r_{ab}^T=0,
\label{orthog}
\stp
for any weights, provided the fit uses the same weights.

\paragraph{The target.}  The dissipative function of a moving fluid contracts both indices of the strain
with the comoving metric (\texttt{particle\_metrics.tex}, ``The metric''; \texttt{stochastic.tex}, ``A moving fluid''):
\st
K_a^{\rm res}=V_aT_a\Big[\eta_a\,\tr\big(\ell_aY_a\ell_aY_a\big)+\tfrac12\zeta_{\rm eff}\,\theta_a^2\Big],\qquad
Y_a=\tfrac12(G_a+G_a^T)-\theta_aP_a,\qquad \theta_a=\frac{\ell_a{:}G_a}{\langle P\ell\rangle_a},
\label{target}
\stp
with $P$ and $\langle P\ell\rangle$ the scalar projector of \texttt{stochastic.tex} and $T_a=T_{\rm rel,a}$.  The question is where the
two factors of $\ell_a$ come from.

%=============================================================================
\section{The stochastic step}
\label{sec:stoch}

\paragraph{The stress move.}  Particle $a$ draws a random symmetric stress $\xi_a$ and gives the traction $\xi_a\bm A_{ab}$
to each bond.  The entropy it produces is, to first order,
\st
\Delta\mathbb S=-\sum_{b\neq a}\Delta\beta_{ab}\cdot\xi_a\bm A_{ab}=-\xi_a{:}\sum_{b\neq a}\Delta\beta_{ab}\bm A_{ab}^T=-V_a\,\xi_a{:}G_a ,
\label{dSstress}
\stp
by \Eq{fit}.  Only the closure $\sum_b\bm A_{ab}\bm r_{ab}^T=V_a\one$ is used, and it holds for both \Eq{Aold} and \Eq{wnew}.  The
stress is drawn as
\st
\xi_a=\sqrt{\frac{2\Delta t\,T_a}{V_a}}\sum_{k=1}^3z_kB_k,\qquad
B_{1,2}=\sqrt{2\eta_a}\,E(e_{1,2}),\quad B_3=\sqrt{\zeta_{\rm eff}}\,\frac{\ell_a}{\langle P\ell\rangle_a},\qquad
E(e)=\ell_a^{1/2}e\,\ell_a^{1/2}-\frac{P{:}\ell_a^{1/2}e\,\ell_a^{1/2}}{\langle P\ell\rangle_a}\,\ell_a ,
\label{B}
\stp
$z_k$ standard normal, $e_{1,2}$ the traceless symmetric unit matrices.  The dissipation of the move is
\st
K_a^{\rm res}=\frac{V_aT_a}{2}\sum_k(B_k{:}G_a)^2=V_aT_a\Big[\eta_a\tr(\ell_aY_a\ell_aY_a)+\tfrac12\zeta_{\rm eff}\theta_a^2\Big],
\label{Kfromxi}
\stp
\Eq{target}.  Indeed $E(e){:}G=e{:}Z$ with $Z=\ell^{1/2}Y\ell^{1/2}$, which is traceless because $\ell{:}Y=\ell{:}G-\theta\,P{:}\ell=0$
($\langle P\ell\rangle=P{:}\ell$), so $\sum_{k=1,2}(E(e_k){:}G)^2=\tr Z^2=\tr(\ell Y\ell Y)$; checked numerically for a random $G$ at
$\bm v=(0.7,0.3)$, to roundoff.  The two factors of $\ell_a$ are the two factors $\ell_a^{1/2}$ on each side of $e$ in
$E(e)$: they are in the Onsager matrix, not in the bond moment.  So the stress move needs \emph{no change} beyond the
weights and the faces \Eq{wnew}.

\paragraph{The residual move.}  Particle $a$ draws $\bm y_{ab}$ on each bond with $\langle\bm y\bm y^T\rangle\propto c_{ab}\,\ell_a$ and gives
$\bm y_{ab}-(Y_a/V_a)\bm A_{ab}$, $Y_a=\sum_b\bm y_{ab}\bm r_{ab}^T$.  The projection removes the resolved part by the
closure alone, $\sum_b[\bm y_{ab}-(Y_a/V_a)\bm A_{ab}]\bm r_{ab}^T=0$.  Its dissipation is the residual
\st
K_a^{\rm resid}=V_aT_a\eta_a\sum_{b\neq a}w_{ab}\,\delta_{ab}\cdot\ell_a\,\delta_{ab}\ \ge0 ,
\label{Kresid}
\stp
zero for linear fields; it damps the modes the fit cannot see.  With lab weights it is round in the lab
instead of the rest frame, which only changes how the unresolved modes are damped.

\paragraph{So:} for the stochastic step, one kernel means $\psi_{ab}$ from \Eq{psicommon} at the lab distance, $\Lambda_a$ for the
flat metric, and $\bm A_{ab}=V_aw_{ab}\bm r_{ab}$.  Nothing else changes.

%=============================================================================
\section{The deterministic step}
\label{sec:det}

The present deterministic bond form,
\st
K=\sum_aV_aT_a\eta_a\sum_{b\neq a}w_{ab}\,\mathcal D_{ab}\cdot\ell_a\,\mathcal D_{ab}+\tfrac12\sum_aV_aT_a\zeta_{\rm eff}\theta_a^2,\qquad
\mathcal D_{ab}=\Delta\beta_{ab}-\big(\Omega_a+\theta_aP_a\big)\bm r_{ab} ,
\label{Kold}
\stp
gets one factor $\ell_a$ explicitly (between the components) and the other from the moment: with
$\mathcal D_{ab}=Y_a\bm r_{ab}+\delta_{ab}$ and \Eq{orthog},
\st
\sum_{b\neq a}w_{ab}\,\mathcal D_{ab}\cdot\ell_a\mathcal D_{ab}=\tr\Big(Y_a\ell_aY_a\,\frac{M_a}{n_a}\Big)+\sum_{b\neq a}w_{ab}\,\delta_{ab}\cdot\ell_a\delta_{ab} .
\label{split}
\stp
With \Eq{Mold} the first term is $\tr(\ell Y\ell Y)$; with \Eq{Mnew} it would be $\tr(Y\ell Y)$, one $\ell$ short.  With one kernel
we therefore write $K$ in its split form directly, the resolved part being the target:
\st
\boxed{\;K=\sum_a\Big(K_a^{\rm res}+K_a^{\rm resid}\Big)
=\sum_aV_aT_a\Big[\eta_a\tr\big(\ell_aY_a\ell_aY_a\big)+\tfrac12\zeta_{\rm eff}\theta_a^2+\eta_a\sum_{b\neq a}w_{ab}\,\delta_{ab}\cdot\ell_a\delta_{ab}\Big].\;}
\label{Knew}
\stp
With \Eq{Mold} this is identical to \Eq{Kold}; with \Eq{Mnew} it is the one-kernel form, and it is the mean of
the stochastic step of \Sect{sec:stoch}.

\paragraph{The force} $\bm F_a=-\partial K/\partial\beta_a$.  $K$ depends on $\beta$ through $G_a$ (linear in $\Delta\beta_{ab}$, \Eq{fit}) and through the
$\delta_{ab}$ at fixed $G_a$.  The derivative of $K_a^{\rm resid}$ through $G_a$ is
\st
\frac{\partial K_a^{\rm resid}}{\partial G_a}=-2V_aT_a\eta_a\,\ell_a\sum_{b\neq a}w_{ab}\,\delta_{ab}\bm r_{ab}^T=0
\stp
by \Eq{orthog}, so the residual acts only through bond tractions,
\st
\bm t_{ab}=-2V_aT_a\eta_a\,w_{ab}\,\ell_a\,\delta_{ab}\qquad(\text{momentum }a\text{ gives }b),
\label{tresid}
\stp
and the resolved part only through the fit, with
\st
E_a\equiv\frac{\partial K_a^{\rm res}}{\partial G_a}
=2V_aT_a\eta_a\Big[\ell_aY_a\ell_a-\frac{P{:}\ell_aY_a\ell_a}{\langle P\ell\rangle_a}\,\ell_a\Big]
+V_aT_a\zeta_{\rm eff}\,\theta_a\,\frac{\ell_a}{\langle P\ell\rangle_a}
\label{E}
\stp
(from $dK^{\rm res}=2V_aT_a\eta_a\,\ell Y\ell{:}dY+V_aT_a\zeta_{\rm eff}\theta\,d\theta$, $dY={\rm sym}\,dG-P\,\ell{:}dG/\langle P\ell\rangle$).  Since
$\partial G_a/\partial\beta_b=w_{ab}\,(\cdot)\,\bm r_{ab}^T(M_a/n_a)^{-1}$ for $b\neq a$ and minus their sum for $b=a$, the resolved force is
again a sum over bonds, $\partial K/\partial\beta_b\mathrel{+}=E_a(M_a/n_a)^{-1}\bm r_{ab}\,w_{ab}$ and the opposite to $a$.  Both parts
are sums of antisymmetric pair forces: momentum is conserved.  $E_a$ is symmetric, so the resolved
part exerts no torque; the residual tractions, as in the present scheme, need not be central.  The
entropy produced in a step $\Delta t$ is $2K\Delta t\ge0$.

\paragraph{The rest frame.}  At $\bm v_a=0$, $\ell_a=\one$, $Y_a=\mathring S_a$ and \Eq{Knew} is the split form (\texttt{bond\_dissipation.tex}, eq.~split2d) of the present \Eq{Kold}: identical.

%=============================================================================
\section{The residual in a moving fluid}
\label{sec:resid}

The resolved part is exact.  The residual \Eq{Kresid} is not: it is a regulator, of relative size $(kh)^2$
for smooth fields, and it knows only the metric of its weights.  For a smooth field the residual of a
bond is the curvature, $\delta_{ab}\simeq\frac12\partial_i\partial_j\beta\,r^i_{ab}r^j_{ab}$, so it is weighted by the fourth moment of the
stencil.  Take the shear wave along the flow, $\beta_y(x)$, $\bm k\parallel\bm v$.  The resolved part is proportional to
$\ell_{xx}=1/\gamma^2$ (time dilation slows the physical rate); the residual is proportional to $\sum_bw_{ab}r_x^4$, which
for comoving weights carries $\ell_{xx}^2$ but for lab weights carries nothing.  The relative error of the
one-kernel residual is therefore enhanced by about $\gamma^2$ (to $\gamma^4$) over the present one, exactly for
the mode of the boost benchmark.

The cure is again to put $\ell$ in by hand: weight the residual of each bond by
\st
f_{ab}=\big(\hat{\bm r}_{ab}\cdot\ell_a\,\hat{\bm r}_{ab}\big)^2 ,
\qquad K_a^{\rm resid}=V_aT_a\eta_a\sum_{b\neq a}w_{ab}\,f_{ab}\,\delta_{ab}\cdot\ell_a\,\delta_{ab},
\label{fab}
\stp
which is $\ell_{xx}^2$ for bonds along the flow and 1 across it, as for the comoving stencil.  $K^{\rm resid}$ stays
non-negative and zero for linear fields.  Two consequences:
\begin{itemize}
\item \emph{Stochastic step:} the residual move couples its noise only to $\delta_{ab}$ (its entropy change is
  $-\sum_b\bm y_{ab}\cdot\delta_{ab}$, by the closure), so any positive coefficient may be used: $c_{ab}\to c_{ab}f_{ab}$.
\item \emph{Deterministic step:} with $f\neq1$ the orthogonality \Eq{orthog} no longer removes the derivative of the
  residual through the fit; it adds $-\sum_bc_{ab}f_{ab}\,\ell_a\delta_{ab}\bm r_{ab}^T$ to $E_a$ \Eq{E}.
\end{itemize}
Both variants are options: \texttt{SPH\_DISSIPATION=common} (with $f_{ab}$) and \texttt{common\_lab} ($f_{ab}=1$).

%=============================================================================
\section{Implementation}
\label{sec:impl}

\begin{itemize}
\item \texttt{SimParams.dissipation = :common} or \texttt{:common\_lab}, set by \texttt{SPH\_DISSIPATION} (recorded in \texttt{[sim]} and
  \texttt{[settings]} of the run record); 2D, with \texttt{hamiltonian = :energy}.  The default stays \texttt{:particle}.
\item Weights (\texttt{src/viscous.jl}): $\psi_{ab}=-W'(r_{ab};h)/r_{ab}$ of \texttt{params.kernel} at the lab distance; $\Lambda_a$ with the flat
  metric; $n_a=\frac12\tr M_a$.
\item Stochastic step (\texttt{build\_stars}): these weights, $\bm A_{ab}=V_aw_{ab}\bm r_{ab}$, and $c_{ab}f_{ab}$; the rest unchanged.
\item Deterministic step (\texttt{\_bond\_dissipation!}): \Eq{Knew} with \Eq{fab}, the tractions \Eq{tresid} (times $f_{ab}$) and the
  mean-field term \Eq{E} plus the residual term above.
\item The present scheme is untouched (every change is behind \texttt{dissipation} $\in\{$\texttt{:common}, \texttt{:common\_lab}$\}$).
\end{itemize}

%=============================================================================
\section{Tests}
\label{sec:tests}

\paragraph{1. At rest} (\texttt{test/testcommonkernel.jl}, C2): $K$ and the forces of the one-kernel step equal those of the
present step to $2\times10^{-16}$ and $1\times10^{-15}$, on a perfect and on a jittered ($0.3a$) lattice; $\sum_a\bm F_a=0$ to roundoff.

\paragraph{2. One mode in a moving fluid} ($v_0=0.8$, $L=1.2$, $h/a=2$): $K$ of $\beta=A\,\bm e\sin\bm k\cdot\bm x$ against the target
\Eq{target}:
\begin{center}
\begin{tabular}{lccc}
\hline
$K/K_{\rm target}$ at $kh=0.26$, $0.52$, $1.05$ & present & \texttt{common\_lab} & \texttt{common} \\ \hline
shear, $\bm k\parallel\bm v$ & $1.003$, $1.013$, $1.049$ & $1.059$, $1.229$, $1.799$ & $0.999$, $0.997$, $0.991$ \\
long., $\bm k\parallel\bm v$ & $0.997$, $0.991$, $0.967$ & $1.017$, $1.067$, $1.229$ & $0.989$, $0.959$, $0.854$ \\
shear, $\bm k\perp\bm v$ & $1.008$, $1.034$, $1.116$ & $1.008$, $1.034$, $1.116$ & $1.003$, $1.014$, $1.040$ \\
long., $\bm k\perp\bm v$ & $1.007$, $1.027$, $1.092$ & $1.007$, $1.027$, $1.091$ & $1.002$, $1.008$, $1.021$ \\
\hline
\end{tabular}
\end{center}
All converge to the target as $(kh)^2$.  With the lab residual the shear wave along the flow is $23\%$ high at
$kh=0.52$ (\Sect{sec:resid}); with $f_{ab}$ it is within $0.3\%$, better than the present scheme, and so are both
modes across the flow.  The price is the longitudinal wave along the flow, $4\%$ low at $kh=0.52$ (the
present scheme: $1\%$), where the weaker residual no longer offsets the low SPH gradient.

\paragraph{3. The viscous step alone,} deterministic (\texttt{examples/commonkernel\_waves.jl}, job 61227129): one mode
$\delta\bm m_a=\epsilon\,\bm e\sin\bm k\cdot\bm x_a$, positions fixed, no ideal step, $v_0=0$ and $0.8$.  The entropy deficit
$D=\sum_a[S_a(e_a,\bar{\bm m})-S_a(e_a,\bm m_a)]$ obeys $\dot D=-2K$, so the mode decays at $\Gamma=K/D$.  In every case the fitted
decay equals $K/D$ of the same scheme to $0.01$--$0.3\%$, the entropy produced equals $2\!\int\!K\,dt$ to $10^{-4}$, and
$\Delta H/H\sim10^{-13}$: the time stepping is consistent with the dissipative function.  At rest the three schemes are
identical.  At $v_0=0.8$, $\Gamma/\Gamma_{\rm hydro}$ at $kh=0.26$, $0.52$, $1.05$:
\begin{center}
\begin{tabular}{lccc}
\hline
 & present & \texttt{common\_lab} & \texttt{common} \\ \hline
shear, $\bm k\parallel\bm v$ & $1.003$, $1.013$, $1.050$ & $1.059$, $1.230$, $1.802$ & $0.999$, $0.997$, $0.992$ \\
long., $\bm k\perp\bm v$ & $1.007$, $1.028$, $1.095$ & $1.007$, $1.028$, $1.095$ & $1.002$, $1.009$, $1.024$ \\
\hline
\end{tabular}
\end{center}
the same as the $K$ ratios of test 2.  (A momentum perturbation \emph{along} a moving flow is not a pure mode of the
viscous step alone: the step holds each particle's energy and $\delta e=\bm v\cdot\delta\bm m$; those modes are covered by test 2.)

\paragraph{4. The stochastic step alone} (positions frozen, \texttt{SR\_NOHAM=1}; \texttt{nersc/campaign\_commonkernel.txt}, job 61227130;
$L=1.2$, 8 runs each, $t=50$):
\begin{center}
\begin{tabular}{lccc}
\hline
 & acceptance, stress & residual & KE$/2VwT$ \\ \hline
present, $v_0=0$ & $0.969$ & $0.773$ & $0.997$ \\
\texttt{common}, $v_0=0$ & $0.969$ & $0.774$ & $0.998$ \\
present, $v_0=0.8$ & $0.968$ & $0.841$ & $17.660$ \\
\texttt{common}, $v_0=0.8$ & $0.981$ & $0.878$ & $17.661$ \\
\hline
\end{tabular}
\end{center}
Both sample the moving equilibrium (predicted KE $17.6$); $\Delta H/H\sim10^{-10}$.  At rest they are identical.  At $v_0=0.8$
the one-kernel step accepts more: its faces are smaller (the lab stencil is not contracted along $\bm v$, so
$\sum_b|\bm A_{ab}|^2$ is $0.33$ of the present one) and its residual coefficients carry $f_{ab}\le1$ ($\sum_bc_{ab}$ is $0.57$).  Without
the ideal step the long waves relax only by diffusion, and the equal-time correlators are not converged
by $t=50$; they are measured in the full benchmark.

\paragraph{5. The boost benchmark} (\texttt{boost.tex}) with \texttt{SPH\_DISSIPATION=common}: running.

\end{document}
